\documentclass[10pt,a4paper]{amsart}
\usepackage[margin=19mm]{geometry}
\usepackage{amsmath,amssymb,mathtools}
\usepackage[scr=rsfs,cal=euler]{mathalfa}
\usepackage{xfrac}
\usepackage[unicode,hidelinks,bookmarks=false]{hyperref}
\newcommand{\ie}{i.e.\ }
\newcommand{\cf}{cf.\ }
\newcommand{\resp}{respectively\ }
\newcommand{\Subsection}{\subsection}
\providecommand{\nobreakdash}{\nobreak\hbox{-}}
\setlength{\emergencystretch}{2em}
\multlinegap0pt
\usepackage{bm}
\usepackage{bbm}
\usepackage{dsfont}
\usepackage{xspace}
\usepackage{cancel}
\usepackage{mathtools}
\usepackage{amsthm}
\makeatletter
\@ifundefined{theorem}{\newtheorem{theorem}{Theorem}}{}
\makeatother
\title[Explicit Runge-Kutta Methods with Multiquadric and Inverse M]{Explicit Runge-Kutta Methods with Multiquadric and Inverse Multiquadric Radial Basis Functions}
\author{Shipra Mahata, Samala Rathan}
\date{Source: arXiv:2507.14474v2 (CC BY 4.0)}
\begin{document}
\maketitle

\noindent\emph{Source and licence.} Explicit Runge-Kutta Methods with Multiquadric and Inverse Multiquadric Radial Basis Functions. Version arXiv:2507.14474v2 (CC BY 4.0), distributed under the Creative Commons Attribution 4.0 International licence, as recorded in the arXiv licence metadata for this exact version and shown on the abstract page https://arxiv.org/abs/2507.14474v2. Full rights notice: licenses/arxiv-2507.14474-CC-BY-4.0.txt.\par\smallskip

\noindent\textit{Editorial scope.} Counted unit: the Theorem whose source label is \texttt{None} in the article (source0.tex lines 896--900, source label \texttt{None}), delivered together with the article's own complete original proof of it (source0.tex lines 901--995, one \texttt{proof} environment of 95 lines). No mathematics is written, repaired or re-derived: statement and proof are the article's own text line for line. Required context retained uncounted: eq:sumb (lines 193--195); eq:sumbn (lines 197--199); eq:rk\_3\_con (lines 279--281); eq:imq3 (lines 563--572). Every definition, display and label of those lines is retained unchanged. Omitted from this package and not redistributed: 1--192, 196--196, 200--278, 282--562, 573--895, 996--1397. Nothing inside a retained excerpt is omitted or altered: every retained body is delivered exactly as the article's lines are written, so no omission receipt of any kind accompanies this package. Citations: the retained excerpts carry no citation command at all, so no bibliography entry of the article is transcribed and no reference is redistributed. Reference adaptation: none was needed and none was made; every reference inside the delivered unit resolves inside it, so no unresolved reference or citation is produced. Printed slips kept literal: no printed word, symbol, hypothesis or number of the counted statement or of its proof was altered. Layout and class: the article's own class is \texttt{elsarticle} with packages including geometry, inputenc, fontenc, amsmath, amssymb, amsthm, breqn, pdflscape, bm, mathrsfs, euscript, oldgerm, enumitem, graphicx; the wrapper is the lane's pinned \texttt{amsart} editorial wrapper at 10pt on A4, which restates the theorem-like environments the retained text needs and adds the article's own macro definitions that the retained text needs (none), with extra wrapper packages (bm, bbm, dsfont, xspace, cancel, mathtools). The delivered numbering is the unit's own. Assets: no retained span contains a figure environment, a graphics-inclusion command, a PDF-page inclusion, a table or a TikZ picture; no image file is copied into this package, the package declares no asset dependency and no image file is redistributed with it. Licence: version arXiv:2507.14474v2 of this article is distributed under the Creative Commons Attribution 4.0 International licence, as recorded in the arXiv licence metadata for this exact version and shown on the abstract page https://arxiv.org/abs/2507.14474v2. A full rights notice is delivered at licenses/arxiv-2507.14474-CC-BY-4.0.txt. Characters: every retained excerpt is pure ASCII, so the delivered engine, XeLaTeX with its default Latin Modern text font, reports no missing character.
\begin{equation}    \label{eq:sumb}
\sum_{j=1}^{r}{b_j} = 1,
\end{equation}

\begin{equation}    \label{eq:sumbn}
\sum_{j=1}^{i-1}{a_{ij}} = c_i, i = 1,...,r.
\end{equation}

\begin{equation}\label{eq:rk_3_con}
b_2 c_2 + b_3 c_3 = \frac{1}{2},\,\,\, b_2 c_2^2 + b_3 c_3^2 = \frac{1}{3},\,\, a_{32} b_3 c_2 = \frac{1}{6},
\end{equation}

\begin{eqnarray}    \label{eq:imq3}  
\begin{cases}
\begin{aligned}
v_{n+1} & = v_n + h(b_1 k_1 + b_2 k_2 + b_3 k_3),\\
k_1 &= f(t_n, v_n), \\
k_2 &= f\left(t_n + c_2 h, \sqrt{(1 + \epsilon_{n2}^2 c_{2}^2 h^2)}ha_{21} f_n + \frac{v_n} {\sqrt{1 + \epsilon_{n2}^2 c_{2}^2 h^2}}\right),\\
k_3 &= f\left(t_n + c_3 h, \sqrt{(1 + \epsilon_{n3}^2 c_{3}^2 h^2)}h (a_{31}k_1 + a_{32}k_2)+ \frac{v_n} {\sqrt{1 + \epsilon_{n3}^2 c_{3}^2 h^2}}\right). 
\end{aligned}    
\end{cases}    
\end{eqnarray}

\begin{theorem} If $\epsilon_{n2}^2$ and $\epsilon_{n3}^2$
are bounded for all $n = 0, 1, \dots, N - 1$. Then the
three-stage RBF Runge-Kutta method (\ref{eq:imq3}), converges provided the method satisfies (\ref{eq:sumb}), (\ref{eq:sumbn}) and
(\ref{eq:rk_3_con}).
\end{theorem}

\begin{proof} Suppose, $E_{n+1}$ be the error $E_{n+1}$  between $u_{n+1}$ and $v_{n+1}$ for fix $t = t_n$. Now,
\begin{align*}
k_1(w) &= f(t_n, w), \\
k_2(w) &= f\left(t_n + c_2h, \sqrt{(1 + \epsilon_{n2}^2 a_{21}^2 h^2)} \,  h a_{21}k_1 + \frac{w} {\sqrt{1 + \epsilon_{n2}^2 a_{21}^2 h^2}}\right), \\
k_3(w) &= f\left(t_n + c_3h, \sqrt{(1 + \epsilon_{n3}^2 c_{3}^2 h^2)} \, h (a_{31}k_1 + a_{32}k_2)+\frac{w} {\sqrt{1 + \epsilon_{n3}^2 c_3^2 h^2}}\right).
\end{align*}
$u_{n+1}$ can be expressed as
\begin{align*}
 u_{n+1} = u_n + h (b_1k_1(u_n) + b_2k_2(u_n) + b_3k_3(u_n)) + h \tau_n,
\end{align*}
where, $\tau_n$ is trancation error.
 Since $f(t,u)$ is Lipschitz continuous in u , so
 \begin{equation*}
     |k_1(u_n) - k_1(v_n)|\leq L|u_n - v_n|,
 \end{equation*}

\begin{equation*}
     |k_2(u_n) - k_2(v_n)|\leq L\left|\sqrt{1 + \epsilon_{n2}^2a_{21}^2h^2}\,\,hk_1a_{21} +\frac{1}{\sqrt{1 +\epsilon_{n2}^2a_{21}^2h^2}}\right||u_n - v_n|,
 \end{equation*}
 \begin{align*}
|k_3(u_n) - k_3(v_n)| 
&= \Bigg| 
f\left(t_n + c_3 h,\ \sqrt{1 + \epsilon_{n3}^2 c_3^2 h^2} \,\, h (a_{31} k_1(u_n) + a_{32} k_2(u_n)) 
+ \frac{u_n}{\sqrt{1 + \epsilon_{n3}^2 c_3^2 h^2}} \right) \\
&\qquad - 
f\left(t_n + c_3 h,\ \sqrt{1 + \epsilon_{n3}^2 c_3^2 h^2} \,\,  h (a_{31} k_1(v_n) + a_{32} k_2(v_n)) 
+ \frac{v_n}{\sqrt{1 + \epsilon_{n3}^2 c_3^2 h^2}} \right)
\Bigg|, \\
&\leq L \Bigg| 
\sqrt{1 + \epsilon_{n3}^2 c_3^2 h^2}  h \left( a_{31}(k_1(u_n) - k_1(v_n)) 
+ a_{32}(k_2(u_n) - k_2(v_n)) \right) 
+ \frac{u_n - v_n}{\sqrt{1 + \epsilon_{n3}^2 c_3^2 h^2}} 
\Bigg|, \\
&\leq L \left[ 
\sqrt{1 + \epsilon_{n3}^2 c_3^2 h^2} \,\, h a_{31} L |u_n - v_n| 
+ \sqrt{1 + \epsilon_{n2}^2 a_{21}^2 h^2} \,\,  h a_{32} L |u_n - v_n| \right. \\
&\quad \left. +  \frac{L|u_n - v_n|}{\sqrt{1 + \epsilon_{n2}^2 a_{21}^2 h^2}}
+ \frac{|u_n - v_n|}{\sqrt{1 + \epsilon_{n3}^2 c_3^2 h^2}} 
\right],
\end{align*}
where, $L$ is Lipschitz constant. 
Now,
\begin{align*}
    E_{n+1} & = |u_{n+1} - v_{n+1}|, \\
& = |(u_n - v_n) + h \left[ b_1(k_1(u_n) - k_1(v_n)) + b_2(k_2(u_n) - k_2(v_n))\right ] + h\tau_n|, \\
& \leq |u_n - v_n| + hb_1L|u_n - v_n| + b_2hL\left(\sqrt{1 + \epsilon_{n2}^2a_{21}^2h^2}\,\,hLa_{21} +  \frac{1}{\sqrt{1 +\epsilon_{n2}^2a_{21}^2h^2}}\right)|u_n - v_n| \\
& + b_3hL \Biggr[ \sqrt{1 + \epsilon_{n3}^2h^2c_3^2}\,\,ha_{31}L  + \sqrt{1+\epsilon_{n3}h^2c_3^2}\,\,ha_{32} \Biggr(L^2\sqrt{1 + \epsilon_{n2}^2a_{21}^2h^2}\,\,ha_{21} + \frac{L}{\sqrt{1 + \epsilon_{n2}^2a_{21}^2h^2}}\Biggl)\Biggr]|u_n - v_n|. 
\end{align*}
Now,
\begin{align*}
    \phi_n &= 1 + b_1 h L 
    + h b_2 L \left( \sqrt{1 + \epsilon_{n2}^2 a_{21}^2 h^2}\,\, h L a_{21} + \frac{1}{\sqrt{1 + \epsilon_{n2}^2 h^2 a_{21}^2}} \right) \\
    &\quad + b_3 h L \left[ \sqrt{1 + \epsilon_{n3}^2 h^2 c_3^2}\,\, h a_{31} L 
    + \sqrt{1 + \epsilon_{n3}^2 h^2 c_3^2}\,\, h a_{32} \left( L^2 \sqrt{1 + \epsilon_{n2}^2 a_{21}^2 h^2}\,\, h a_{21} 
    + \frac{L}{\sqrt{1 + \epsilon_{n2}^2 a_{21}^2 h^2}} \right) \right], \\
    & = 1 + hL \left( b_1 + \frac{b_2}{\sqrt{1 + \epsilon_{n2}^2h^2a_{21}^2}} + \frac{b_3}{\sqrt{1 
     +\epsilon_{n3}^2h^2c_3^2}}\right) + h^2L^2\left(b_2a_{21}\sqrt{1 + \epsilon_{n2}^2a_{21}^2h^2} +  b_3a_{31}\sqrt{1 + \epsilon_{n3}^2h^2c_3^2} \right. \\
     & \left. + a_{32}b_3\sqrt{1 + \epsilon_{n3}^2h^2c_3^2}\sqrt{1 + \epsilon_{n2}^2a_{21}^2h^2}\right) + b_3L^3h^3a_{21}a_{32}\sqrt{1 + \epsilon_{n3}^2h^2c_3^2}\sqrt{1 + \epsilon_{n2}^2h^2a_{21}^2}, \\
     & \leq 1 + hL + \frac{h^2L^2}{2} - \frac{h^2L^2}{2} \left( b_2a_{21}\sqrt{1 + \epsilon_{n2}^2a_{21}^2h^2} +  b_3a_{31}\sqrt{1 + \epsilon_{n3}^2h^2c_3^2} \right. \\
     & \left. + a_{32}b_3\sqrt{1 + \epsilon_{n3}^2h^2c_3^2} \right) + b_3L^3h^3a_{21}a_{32}\sqrt{1 + \epsilon_{n3}^2h^2c_3^2}\sqrt{1 + \epsilon_{n2}^2h^2a_{21}^2}, \\
      & \leq (1 + \mathcal{C}_3 h) e^{hL}, \\
      & \leq e^{h(1+\mathcal{C}_3)},
\end{align*}
where,
 \begin{align*}
    \mathcal{C}_3 = \sup_{\substack{0 < h \leq b - a \\ n = 0, \dots, N - 1}} 
    \frac{A}{ e^{hL} },
\end{align*}
with 
\begin{align*}
A=  & h^2 L^3 b_3 c_{2} a_{32} \sqrt{1 + \epsilon_{n3}^2 c_3^2 h^2} \sqrt{1 + \epsilon_{n2}^2 c_{2}^2 h^2} \\
&\qquad- \frac{h^2L^2}{2} \biggl( b_2 c_{2} \sqrt{1+\epsilon_{n2}^2 c_{2}^2 h^2}+ b_3 a_{31} \sqrt{1 + \epsilon_{n3}^2 h^2 c_3^2} 
          + b_3 a_{32} \sqrt{1+ \epsilon_{n3}^2 c_3^2 h^2} \biggr).
\end{align*}
Here, $b_3$, $c_2$, $c_3$, $a_{32}$ is constant and $\epsilon_{n2}^2$ and $\epsilon_{n3}^2$ is bounded as it depends on $u$ and $u$ is bounded ensured by Lipschitz continuity. Therefore the term $E_n$ satisfies 
\begin{equation*}
    E_n \leq \sum_{j=0}^{n-1} \varphi_j E_0 + h \sum_{j=1}^{n-1} \left( \sum_{m=j}^{n-1} \varphi_m \right) |\tau_{j-1}| + h |\tau_{n-1}|.
\end{equation*}
Clearly,
\begin{equation*}
    \sum_{m=j}^{n-1} \varphi_m \leq \sum_{m=0}^{n-1} \varphi_m \leq e^{h(L + C_2) n} \leq e^{h(L + C_2) N} = e^{(L + C_2)(b - a)}.
\end{equation*}
With \( E_0 = 0 \) and
\[
\|\tau\|_\infty = \max_{n = 0, \dots, N-1} |\tau_n|,
\]
it can be concluded that for every \( n = 0, 1, \dots, N \),
\begin{equation*}
    E_n \leq h e^{(L + C_2) T} \sum_{j=0}^{n-1} \|\tau\|_\infty = n h e^{(L + C_2)(b - a)} \|\tau\|_\infty \leq (b - a) e^{(L + C_2)(b - a)} \|\tau\|_\infty.
\end{equation*}
Hence,
\begin{equation*}
\lim_{\substack{h \to 0 \\ N h = b - a}} E_n = 0.
\end{equation*}
\end{proof}

\end{document}
